\section{Conclusion}

This paper presented PC-RCO, a CFD-Prior Current Observer that combines continuous Gauss-Markov regularization with adaptive update-rate clamping for robust AUV current estimation under hydrodynamic model mismatch. The main contributions and findings are:

\begin{enumerate}[itemsep=0pt,topsep=4pt]
    \item \textbf{Model mismatch robustness}: Under 30\% CFD drag perturbation, PC-RCO degrades by only 42\% versus 77\% for a CFD-augmented EKF, with the robustness advantage widening to 88 percentage points at 50\% mismatch. At zero mismatch, PC-RCO achieves 3$\times$ better accuracy than the kinematic baseline (RMSE$_{V_c}=0.064$ vs.\ 0.192).

    \item \textbf{Continuous GM regularization}: Replacing binary excitation gating with continuous GM decay applied at every time step eliminates the need for threshold tuning while providing 57--218\% improvement in estimation accuracy under mismatch and high-noise conditions. We demonstrate that binary excitation gating cannot reliably distinguish excitation levels in feedback-controlled AUVs---a finding with implications for observer design beyond the current work.

    \item \textbf{Adaptive clamping with online noise estimation}: The per-step update bound dynamically scales from 0.06 to 0.18\,m/s based on real-time DVL noise estimates, achieving a 72\% reduction in high-noise RMSE (0.86$\rightarrow$0.24) compared to fixed clamping, while preventing the 29--33$\times$ larger per-step jumps observed in unclamped variants.

    \item \textbf{Multi-objective evaluation framework}: We demonstrate that RMSE alone can be misleading for current observer evaluation---the unclamped variant achieves lower RMSE (0.126 vs.\ 0.243) but produces physically impossible 2.9\,m/s single-step estimation jumps. Per-step update magnitude, overshoot, and final bias should be reported alongside conventional RMSE.
\end{enumerate}

Future work will pursue: (1) water tank and at-sea experimental validation with independent ADCP measurements as current ground truth; (2) integration of PC-RCO into closed-loop MPC to quantify downstream control benefits of bounded current estimates; (3) extension to adaptively update the CFD prior based on accumulated operational data to address long-term model drift from biofouling and structural changes.
